\section{Réalisation d'un GRAFCET de prise/dépose}
\subsection{Un GRAFCET d'approche}
Afin de généraliser la prise et la dépose d'objets par le robot, on propose d'écrire un GRAFCET de \texttt{prise/dépose} qui sera utilisé à chaque prise ou dépose d'un gobelet. Le squelette de ce GRAFCET vous a été donné au tableau et se trouve dans la section \texttt{F1\_EVAL\_PriseDepose} et sa section POST associée.

Pour réaliser ce comportement, on propose d'utiliser les variables de type \texttt{T\_CartesianPos} suivantes
\begin{center}
  \begin{tabular}{l|l|l|p{5cm}}
    Type & Nom                                     & Coordonnées                   & Description                           \\\hline\hline
    T\_CartesianPos & \texttt{cartesianPos\_Cuve0}            & (360,-160,  20, 180,  0,  0)  & Position du gobelet sur la cuve 0     \\
    T\_CartesianPos & \texttt{cartesianPos\_Cuve1}            & (450,   0,  20, 180,  0,  0)  & Position du gobelet sur la cuve 1     \\
    T\_CartesianPos & \texttt{cartesianPos\_Convoyeur}        & ( 90,-410,  50, 180,  0,  0)  & Position du gobelet sur le convoyeur  \\\hline
    T\_Trsf & \texttt{Trsf\_Approche}                 & (  0,   0, 100,   0,  0,  0)  & Translation d'approche                \\\hline
    T\_CartesianPos & \texttt{cartesianPos\_Destination}      &                               & Variable d'échange entre GRAFCET principal et GRAFCET de mouvements. \\
  \end{tabular}
\end{center}

Chaque bloc fonctionnel UniVALplc de mouvement est associé à une variable structurée reliée à ses entrées et sorties. L'appel des bloc est fait dans la section \texttt{ScaraMoves} qu'il est possible de consulter.

A chaque fois que l'on souhaite faire une prise ou dépose de gobelet, la stratégie proposée, pour l'étape correspondante, est la suivante :
\begin{enumerate}
  \item Dans le programme principale \texttt{F1\_ProductionNormale}, une action \textbf{P1} est lancée
    \begin{itemize}
      \item Elle définit la variable \texttt{cartesianPos\_Destination} avec la position de destination voulue. \textit{On utilisera ici les positions définies dans les variables définies ci-dessus.}
      \item Elle positionne une variable \texttt{priseDepose} qui indiquera si le robot doit poser ou récupérer un gobelet.
    \end{itemize}
  \item Le GRAFCET coopérant de la section \texttt{F1\_EVAL\_PriseDepose} se déclenche alors \textit{(Condition de sortie de l'étape initiale à définir)}
    \begin{enumerate}
      \item Le point d'approche est calculé à partir du point de destination
      \item Le robot est envoyé au point d'approche
      \item Le robot descend jusqu'au point de destination
      \item La prise ou dépose est effectuée
      \item Le robot remonte
      \item Le GRAFCET attend que l'étape du programme principal correspondante ne soit plus active
    \end{enumerate}
  \item Le programme principal se poursuit.
\end{enumerate}

\begin{UPSTIManipulation}{Réalisation du GRAFCET de prise/dépose}

  \UPSTIetape{Vérifier, et définir si besoin, les variables du tableau précédent.}

  Ouvrir la documentation afin de prendre connaissance du fonctionnement du bloc \texttt{VAL\_ShiftPoint}.
  \UPSTIetape{Dans la section \texttt{F1\_EVAL\_PriseDepose\_Post}, Appeler une instance de bloc \texttt{VAL\_ShiftPoint}.}
  \UPSTIetape{Dans la section \texttt{F1\_EVAL\_PriseDepose\_Post}, utiliser ce bloc fonctionnel pour faire le calcul du point d'approche lors de l'étape \texttt{F1\_M\_01}}
  \UPSTIetape{Provoquer le lancement du mouvement du robot pour aller au point d'approche. \\\textbf{Attention, il peut être judicieux de vérifier que le robot ne se trouve pas déjà au point d'approche avant de lancer ce mouvement.}}
  \UPSTIetape{Compléter le reste du programme pour suivre le comportement désiré}
  \UPSTIetape{Faire \textbf{Valider} par l'enseignant}
\end{UPSTIManipulation}